\documentclass[11pt]{article}
\usepackage[a4paper,margin=2.3cm]{geometry}
\usepackage{xcolor}
\usepackage{parskip}
\usepackage{titlesec}
\usepackage{enumitem}
\usepackage[most]{tcolorbox}
\usepackage{fancyhdr}

\definecolor{accent}{RGB}{0,51,102}
\definecolor{part1color}{RGB}{0,110,50}
\definecolor{part2color}{RGB}{180,30,30}
\definecolor{softbg}{RGB}{245,245,250}
\definecolor{pitchbg}{RGB}{240,247,255}

\titleformat{\section}{\Large\bfseries\color{accent}}{\thesection.}{0.5em}{}
\titleformat{\subsection}[hang]{\large\bfseries\color{accent}}{}{0em}{}

\pagestyle{fancy}
\fancyhf{}
\lhead{\textit{Thesis Presentation Pitch --- Sanskruti Kolgane}}
\rhead{\thepage}
\renewcommand{\headrulewidth}{0.3pt}

\newtcolorbox{pitchbox}[1]{
  colback=pitchbg,
  colframe=accent!60,
  boxrule=0.3pt,
  arc=2pt,
  title=#1,
  fonttitle=\bfseries\small,
  coltitle=white,
  colbacktitle=accent,
  attach boxed title to top left={xshift=6pt,yshift=-3pt},
  boxed title style={sharp corners, rounded corners=north, boxrule=0pt},
  top=10pt
}

\newtcolorbox{tipbox}{
  colback=yellow!10,
  colframe=orange!70,
  boxrule=0.3pt,
  arc=2pt,
  left=6pt,right=6pt,top=4pt,bottom=4pt
}

\title{\vspace{-1cm}\Huge\textbf{Thesis Presentation Pitch}\\[4pt]
       \large Detailed Speaking Notes in Simple Language --- Slide by Slide}
\author{Kolgane Sanskruti Sanjay \quad (21074015)\\
        M.\,Tech Thesis Viva --- IIT (BHU) Varanasi\\
        \textit{Supervisor:} Dr.\ Anil Kumar Singh}
\date{}

\begin{document}
\maketitle
\thispagestyle{empty}

\vspace{-0.5cm}

\begin{tcolorbox}[colback=softbg,colframe=accent,boxrule=0.5pt]
\textbf{How to use this document.} This is a slide-by-slide speaking guide in simple language. For each slide you will see:
\begin{itemize}[leftmargin=*,itemsep=2pt]
  \item \textbf{What the slide shows} --- a one-line summary.
  \item \textbf{What to say} --- the actual pitch in simple, conversational language.
  \item \textbf{Emphasise} --- the one or two things to stress while speaking.
  \item \textbf{If asked} --- likely follow-up and a short honest answer.
\end{itemize}
\textbf{Narrative spine:} Part 1 is about my novel alignment architecture (Systems F, G, H). Part 2 is about the real Bhojpuri benchmark (Systems A, K). Don't blend the two --- finish Part 1 cleanly before starting Part 2.
\end{tcolorbox}

\vfill
\noindent\textit{Total slides: 26. Target speaking time: 18--22 minutes, leaving 8--12 minutes for questions.}

\clearpage

% ─────────────────────────────────────────────────
\section{Slide 1 --- Title Page}

\begin{pitchbox}{What the slide shows}
Thesis title, author, supervisor, institute.
\end{pitchbox}

\begin{pitchbox}{What to say}
``Good morning, sir. Thank you for being here. My thesis is about building a dependency parser for Bhojpuri, which is a low-resource Indian language. The approach I used is to transfer knowledge from Hindi using parallel bottleneck adapters and a new syntax-aware transfer technique. I'll walk you through everything in two clear parts --- first my novel architecture work, and then the real-world benchmark evaluation.''
\end{pitchbox}

\begin{tipbox}
\textbf{Emphasise:} confident opening, clearly name the two-part structure.
\end{tipbox}

% ─────────────────────────────────────────────────
\section{Slide 2 --- Outline}

\begin{pitchbox}{What the slide shows}
Two coloured blocks listing Part 1 and Part 2 contents.
\end{pitchbox}

\begin{pitchbox}{What to say}
``Before diving in, let me tell you how I've organised this talk. Part 1 is about improving the parser architecture on the Hindi-Bhojpuri aligned data. I have three systems there: System F is a fine-tuning baseline, System G is my first novel contribution using parallel adapters with alignment, and System H is my main architectural contribution --- I call it SACT, syntax-aware cross-lingual transfer.

Part 2 is about the real benchmark, the Bhojpuri Treebank or BHTB. There I have System A as the zero-shot baseline, and System K as my novel UD-Bridge self-training approach. Part 1 builds the architecture and lessons, and Part 2 applies those lessons to the actual benchmark where Bhojpuri parsing is evaluated.''
\end{pitchbox}

\begin{tipbox}
\textbf{Emphasise:} ``Part 1 builds, Part 2 applies'' --- set the listener's mental map upfront.
\end{tipbox}

% ─────────────────────────────────────────────────
\section{Slide 3 --- Dependency Parsing in One Slide}

\begin{pitchbox}{What the slide shows}
Task definition on the left; an example Bhojpuri sentence parsed as a tree on the right.
\end{pitchbox}

\begin{pitchbox}{What to say}
``Let me quickly explain dependency parsing. It's the task of finding the grammatical structure of a sentence. The rule is simple: every word points to exactly one head word, and the arrow between them carries one label that names the grammatical relationship. Together this forms a tree rooted at a special ROOT node.

Look at the example on the right --- `Ram ghar jaa rahaa hai', which means `Ram is going home'. `Ram' is the subject of the verb `jaa', with label nsubj. `Ghar', meaning home, is an oblique argument with label obl. `Rahaa' and `hai' are auxiliary verbs. Each arc has one precise label.

Why does this matter? Because dependency parses feed into nearly every downstream NLP task --- machine translation, question answering, information extraction. You can't build good Bhojpuri NLP tools without a parser.

Two accuracy metrics. UAS, unlabelled attachment score, checks if each word's head is correct, ignoring the label. LAS, labelled attachment score, is stricter --- both head and label must be correct. LAS is the main number I'll report throughout.''
\end{pitchbox}

\begin{tipbox}
\textbf{If asked:} ``Why biaffine?'' $\to$ ``Biaffine scorers have been state-of-the-art for dependency parsing since 2017 (Dozat \& Manning). They model head-dependent interactions with a bilinear form, which is expressive but efficient.''
\end{tipbox}

% ─────────────────────────────────────────────────
\section{Slide 4 --- Why Bhojpuri?}

\begin{pitchbox}{What the slide shows}
Language facts, linguistic leverage, and the available training resources.
\end{pitchbox}

\begin{pitchbox}{What to say}
``Bhojpuri is a large but under-resourced language. It has between 50 and 60 million native speakers --- comparable to Italian or Korean. It's spoken mainly in Bihar, UP and Jharkhand, with a large diaspora in Mauritius, Trinidad, Fiji, Guyana, and Suriname.

Despite its size, Bhojpuri has been historically mislabelled as a dialect of Hindi. It only got its own ISO language code --- `bho' --- in 2006. Because of this underclassification, there's almost no NLP infrastructure for it.

The challenge for my work is there's no large annotated Bhojpuri treebank, so I have to transfer knowledge from Hindi. Fortunately this is linguistically reasonable: Bhojpuri and Hindi share the same SOV word order, a lot of vocabulary, and the same Devanagari script. So transfer is not just possible, it's natural.

I have two resources to work with: a Hindi-Bhojpuri aligned corpus of roughly 30,000 sentence pairs, and the Hindi HDTB treebank with 13,304 annotated sentences in UD format. These two resources are what all my systems use.''
\end{pitchbox}

\begin{tipbox}
\textbf{If asked:} ``Why not use direct Bhojpuri treebanks?'' $\to$ ``Only 357 gold-annotated Bhojpuri sentences exist in UD (the BHTB treebank). That's not enough to train a parser from scratch --- we have to transfer.''
\end{tipbox}

% ─────────────────────────────────────────────────
\section{Slide 5 --- Building Blocks}

\begin{pitchbox}{What the slide shows}
Backbone (XLM-R), Trankit toolkit, Pfeiffer adapters, and the biaffine scorer.
\end{pitchbox}

\begin{pitchbox}{What to say}
``Here are the technical building blocks I use throughout. The backbone is XLM-RoBERTa, or XLM-R. It's a multilingual transformer pretrained on 100 languages --- including both Hindi and Bhojpuri. Because both languages share Devanagari, XLM-R already has sub-word tokens that work across both.

I use the Trankit toolkit, which is an end-to-end multilingual NLP pipeline that supports language-specific adapters and includes a native biaffine parsing head out of the box.

The adapters are key. I use Pfeiffer-style bottleneck adapters with dimension 64. The equation on the slide shows the standard form --- down-projection, non-linearity, up-projection, then residual connection. The important thing is: with adapters, the backbone stays frozen. I only train about 100K parameters per language --- a tiny fraction of the full model. This means no catastrophic forgetting of the multilingual prior.

Finally, the biaffine scorer. It computes the likelihood of each word being the head of each other word using a bilinear form. This is the standard state-of-the-art design for dependency parsing.''
\end{pitchbox}

\begin{tipbox}
\textbf{Emphasise:} frozen backbone + small adapters $\Rightarrow$ efficient transfer without forgetting. This pattern matters for Systems G, H, and K.
\end{tipbox}

\clearpage

% ─────────────────────────────────────────────────
\begin{center}
{\LARGE\bfseries\color{part1color}PART 1 --- Novel Cross-Lingual Alignment}
\end{center}

\vspace{0.3cm}
\begin{tcolorbox}[colback=part1color!5,colframe=part1color,boxrule=0.5pt]
\textbf{Part 1 story in one line:} naive fine-tuning (F) gives poor results on aligned data; my novel adapter-based alignment (G, H) improves it by over \textbf{32 LAS points}, and my syntax-aware variant (H) also converges \textbf{5$\times$ faster}.
\end{tcolorbox}

% ─────────────────────────────────────────────────
\section{Slide 6 --- System F: Warm-Start Fine-Tuning (Baseline)}

\begin{pitchbox}{What the slide shows}
Setup, Dev result (weak), limitations, and the research question that motivates G and H.
\end{pitchbox}

\begin{pitchbox}{What to say}
``Now we start Part 1. System F is my baseline --- the straightforward, obvious approach. I start from a Hindi-pretrained Trankit checkpoint and fine-tune all parameters on the 30,000 Hindi-Bhojpuri aligned pairs using standard cross-entropy loss.

The result is quite weak --- only 27.57\% UAS and 17.75\% LAS on the Dev set. That's surprisingly low, and it tells us something important. A few reasons. First, full fine-tuning updates around 270 million parameters on a small dataset --- that's a recipe for overfitting and catastrophic forgetting of the multilingual knowledge XLM-R already has. Second, there's no explicit cross-lingual alignment signal --- the model isn't told that Hindi and Bhojpuri are related. It just sees Bhojpuri and tries to predict labels.

This motivates my research question: can we do much better by explicitly aligning Hindi and Bhojpuri representations? And the plan is --- System G adds freeze plus MSE alignment, and System H replaces MSE with syntax-aware losses.''
\end{pitchbox}

\begin{tipbox}
\textbf{Emphasise:} F is \textit{deliberately} the naive baseline. The weakness of F makes the Part 1 story dramatic --- my novel systems produce huge gains, not small ones.
\end{tipbox}

% ─────────────────────────────────────────────────
\section{Slide 7 --- System G: Parallel Adapters + MSE Alignment (Novel \#1)}

\begin{pitchbox}{What the slide shows}
Two design changes (freeze backbone, MSE alignment), training objective, and the Dev result.
\end{pitchbox}

\begin{pitchbox}{What to say}
``System G is my first novel contribution. It makes two design changes compared to System F.

First, I freeze XLM-R completely. Instead of fine-tuning the whole 270-million-parameter backbone, I only train small adapter modules plus the parsing heads. Total trainable parameters drop to around 200K --- a thousand times smaller. This prevents catastrophic forgetting because the multilingual prior stays untouched.

Second --- and this is the novel idea --- I add an MSE alignment loss. For matched tokens between Hindi and Bhojpuri in the parallel data, I pull their hidden state representations closer together using mean squared error. The idea is: Hindi has a lot of training signal, Bhojpuri has little. By making Bhojpuri representations mimic Hindi representations at matched positions, I transfer Hindi's structural knowledge into Bhojpuri's representation space.

The full training objective combines the Bhojpuri parsing loss, a Hindi auxiliary loss at weight 0.3, and the MSE alignment at weight 0.5.

The result is 55.70\% UAS and 50.02\% LAS on Dev. That's a massive jump over System F --- \textbf{plus 32.3 LAS points} --- and with a thousand times fewer trainable parameters. The takeaway is: freezing plus explicit alignment completely transforms the performance.''
\end{pitchbox}

\begin{tipbox}
\textbf{Emphasise:} +32.3 LAS gain \textit{and} 1000$\times$ fewer params. Both numbers together tell the story.
\end{tipbox}

% ─────────────────────────────────────────────────
\section{Slide 8 --- System H: SACT (Novel \#2, Main Contribution)}

\begin{pitchbox}{What the slide shows}
Three syntax-aware losses: content-word cosine, Arc-KL distillation, Cross-lingual Tree Supervision (CTS).
\end{pitchbox}

\begin{pitchbox}{What to say}
``System H is my main architectural contribution. I call it SACT --- Syntax-Aware Cross-Lingual Transfer.

The problem with MSE alignment in System G is that it's too blunt. It pulls all representations together with equal force --- including function words, punctuation, and tokens that don't carry syntactic structure. MSE treats every token the same way.

SACT replaces that blunt MSE with three smarter, syntax-aware losses.

\textit{First, content-word cosine alignment.} Instead of MSE, I use cosine similarity --- which is scale-invariant. And I weight tokens by part of speech: nouns and verbs, which carry structure, get weight 1.5. Function words get weight 0.3. This forces the model to focus its alignment effort on syntactic anchors.

\textit{Second, Arc-KL distillation.} I use the Hindi parser as a teacher. For matched tokens, I make the Bhojpuri parser match the Hindi parser's arc score distribution using KL divergence. This transfers not just what the representation looks like --- but what parsing decision Hindi would make.

\textit{Third, Cross-lingual Tree Supervision, CTS.} For each matched Hindi-Bhojpuri token pair, I reuse the Hindi gold head as supervision for Bhojpuri. This doubles my structural training signal using clean gold data --- not auto-transferred data.

The combined objective has five terms with the weights shown. Final Dev result: 55.85\% UAS and 50.08\% LAS --- my best architecture on the aligned data.''
\end{pitchbox}

\begin{tipbox}
\textbf{If asked:} ``Why these three specific losses?'' $\to$ ``Each attacks a different level of supervision: cosine aligns representations, Arc-KL aligns decisions, CTS provides direct structural supervision. The ablation (next slide) shows each is non-redundant.''
\end{tipbox}

% ─────────────────────────────────────────────────
\section{Slide 9 --- Results on Aligned Dev Set}

\begin{pitchbox}{What the slide shows}
Full F/G/H results table plus head-to-head LAS gains.
\end{pitchbox}

\begin{pitchbox}{What to say}
``This slide puts Part 1 together in one table. F gets only 17.75\% LAS. G jumps to 50.02\% --- a plus 32.3 point gain just by freezing the backbone and adding MSE alignment. H reaches 50.08\% --- a tiny 0.06-point gain over G at convergence.

You might look at that 0.06 and think SACT barely matters. But look at the head-to-head column on the right. \textbf{At epoch 1, H is already 3.3 points ahead of G.} So SACT converges dramatically faster --- the syntax-aware losses give stronger gradient signal from the very start. I'll show this visually on the next slide.

The big takeaway, highlighted in green: both of my novel designs --- G and H --- massively beat the naive fine-tuning baseline F.''
\end{pitchbox}

\begin{tipbox}
\textbf{Emphasise:} the 0.06 at convergence is misleading in isolation --- the real win for H is training efficiency.
\end{tipbox}

% ─────────────────────────────────────────────────
\section{Slide 10 --- Learning Curves: G vs.\ H}

\begin{pitchbox}{What the slide shows}
Dev LAS vs.\ epoch for G (MSE) and H (SACT).
\end{pitchbox}

\begin{pitchbox}{What to say}
``This is where SACT really shows off. The x-axis is training epochs, y-axis is Dev LAS.

System G --- the blue squares --- starts low at epoch 1 and climbs slowly over 10 epochs to reach its final score around 50.

System H --- the red triangles --- starts 3.3 points higher at epoch 1 and almost plateaus by epoch 2.

Why does this happen? Because MSE alignment in G only aligns hidden representations. The parser has to then figure out, on its own, how aligned representations should translate into correct parsing decisions. That takes time. SACT's Arc-KL and CTS losses directly supervise parsing decisions --- so the gradient signal is much stronger from epoch 1.

This is practically valuable. If you're compute-constrained or time-constrained, SACT gets you to near-final performance in 2 epochs instead of 10. That's roughly a 5$\times$ speedup.''
\end{pitchbox}

% ─────────────────────────────────────────────────
\section{Slide 11 --- Ablation}

\begin{pitchbox}{What the slide shows}
LAS drop when each SACT component is removed individually, plus removing all auxiliaries.
\end{pitchbox}

\begin{pitchbox}{What to say}
``I ran ablation experiments to confirm every SACT component pulls its weight. I removed one loss at a time and measured the Dev LAS drop.

The ranking, from most important to least: Hindi regularisation loss contributes 0.46 points, cosine alignment contributes 0.37, Cross-lingual Tree Supervision contributes 0.21, Arc-KL distillation contributes 0.14.

Removing all four auxiliary losses drops 1.15 points total --- which is close to the sum of individual drops. This means the four losses are largely complementary, not redundant.

None of the four components is dead weight. Every one of them pulls the architecture forward. This is an important check --- it means the SACT design is principled, not over-engineered.''
\end{pitchbox}

\begin{tipbox}
\textbf{If asked:} ``Arc-KL has lowest LAS drop but highest weight (2.0) in the loss. Why?'' $\to$ ``Arc-KL converges fastest and stabilises training. Its marginal drop at convergence is small because by then the model has learned the decisions, but its role in guiding early training is large --- which is what the learning curves on the previous slide also confirm.''
\end{tipbox}

% ─────────────────────────────────────────────────
\section{Slide 12 --- Part 1 Takeaway (Bridge Slide)}

\begin{pitchbox}{What the slide shows}
A green takeaway box with the Part 1 conclusion, plus a footnote hinting at Part 2.
\end{pitchbox}

\begin{pitchbox}{What to say}
``Let me summarise Part 1 clearly. \textbf{My novel cross-lingual alignment losses improve massively over naive fine-tuning on the Hindi-Bhojpuri aligned corpus.}

What worked? Three things. One --- freezing the backbone and training adapters (G and H) beats full fine-tuning (F) by a huge margin. Two --- syntax-aware supervision in H beats uniform MSE in G on convergence speed. Three --- SACT gives a much stronger gradient signal from epoch 1.

The architectural progression is F to G to H. F to G gives plus 32.3 LAS at convergence. G to H gives only 0.06 at convergence but 3.3 at epoch 1.

[\textit{Pause for a beat.}]

Now --- aligned-data performance is one side of the coin. But the real Bhojpuri benchmark isn't the aligned Dev set. It's something else entirely. Let me turn to that now.''
\end{pitchbox}

\begin{tipbox}
\textbf{Emphasise:} this is the most important bridge slide. Pause after ``let me turn to that now'' for a deliberate beat before moving to Slide 13.
\end{tipbox}

\clearpage

% ─────────────────────────────────────────────────
\begin{center}
{\LARGE\bfseries\color{part2color}PART 2 --- Real-World Benchmark Evaluation}
\end{center}

\vspace{0.3cm}
\begin{tcolorbox}[colback=part2color!5,colframe=part2color,boxrule=0.5pt]
\textbf{Part 2 story in one line:} the actual Bhojpuri benchmark is BHTB, not the aligned Dev set. My Part 1 systems can't be evaluated there because of schema mismatch. System A is the zero-shot baseline (35.36 LAS); my novel System K --- UD-Bridge via silver self-training --- reaches \textbf{36.70 LAS}, beating A by \textbf{+1.34 LAS}.
\end{tcolorbox}

% ─────────────────────────────────────────────────
\section{Slide 13 --- The Bhojpuri Treebank (BHTB) --- The Real Benchmark}

\begin{pitchbox}{What the slide shows}
BHTB facts on the left; the schema mismatch between aligned data and BHTB on the right.
\end{pitchbox}

\begin{pitchbox}{What to say}
``This is the most important motivation slide in Part 2, so please follow carefully.

The Bhojpuri Treebank, or BHTB, is the only gold-annotated Bhojpuri treebank that exists. It has 357 manually annotated sentences in the strict Universal Dependencies schema, and it's the standard evaluation set for any Bhojpuri parser. It's test-only --- we cannot train on it.

Now, you might ask: why can't I just take my Part 1 systems --- F, G, H --- and evaluate them on BHTB? They're good parsers, right? Here's the critical problem.

The Hindi-Bhojpuri aligned corpus I used in Part 1 doesn't have clean UD labels. Its annotations were automatically transferred from Hindi --- and that process introduces multiple roots per sentence, cycles, and relation labels that aren't in the UD inventory. So F, G, and H have all learned this `aligned-data dialect' of labels --- which doesn't match BHTB's strict UD dialect at all.

The label inventories are fundamentally incompatible. If I train in one dialect and evaluate in the other, the training signal lives in the wrong label space. To compete on BHTB, I need a completely different strategy --- I need to train in the UD label space from the start. That's what Part 2 is about.''
\end{pitchbox}

\begin{tipbox}
\textbf{Emphasise:} slow down here. The schema-mismatch explanation is the entire motivation for System K. If the supervisor understands this slide, Part 2 lands easily.
\end{tipbox}

% ─────────────────────────────────────────────────
\section{Slide 14 --- System A: Zero-Shot Baseline on BHTB}

\begin{pitchbox}{What the slide shows}
Zero-shot pipeline diagram on the right; result and baseline role on the left.
\end{pitchbox}

\begin{pitchbox}{What to say}
``System A is my baseline for the BHTB evaluation. The idea is simple: rely entirely on XLM-R's multilingual shared representations.

I train Trankit only on Hindi HDTB --- 13,000 UD-labelled Hindi sentences --- with zero Bhojpuri training. Then I apply it directly to Bhojpuri at test time. Because HDTB is native UD, the model's labels automatically live in the UD schema --- which matches BHTB.

Result: 52.78\% UAS and 35.36\% LAS on BHTB gold. A reasonable baseline considering I haven't seen a single labelled Bhojpuri sentence. It works because XLM-R's shared Hindi-Bhojpuri sub-word space transfers quite naturally.

This 35.36\% LAS is the target --- this is what System K has to beat.''
\end{pitchbox}

\begin{tipbox}
\textbf{If asked:} ``Why does zero-shot work at all?'' $\to$ ``XLM-R is pretrained on both Hindi and Bhojpuri, so its sub-word embeddings already encode cross-language similarity. Combined with the high structural similarity between Hindi and Bhojpuri, the parser generalises --- especially for morphologically-marked local relations.''
\end{tipbox}

% ─────────────────────────────────────────────────
\section{Slide 15 --- System K: UD-Bridge via Silver Self-Training (Novel)}

\begin{pitchbox}{What the slide shows}
Three-step pipeline: generate silver with A $\to$ merge with HDTB $\to$ train K.
\end{pitchbox}

\begin{pitchbox}{What to say}
``System K is my novel contribution for the benchmark. I call it UD-Bridge via silver self-training.

The premise is simple but powerful: \textbf{schema consistency beats data quantity}. What I mean is --- having noisy data in the right label space is more useful than having clean data in the wrong label space.

The pipeline has three steps, shown in the diagram.

\textit{Step 1.} Take the 30,966 Bhojpuri sentences from the Hindi--Bhojpuri aligned corpus --- the same Bhojpuri text used in Part 1 --- and \textbf{strip away the auto-transferred labels}. Run System A on the bare tokens. Because System A operates in UD space, it produces silver UD-labelled Bhojpuri data. The labels are noisy because System A isn't perfect, but crucially they are schema-consistent with BHTB.

\textit{Step 2.} Concatenate this silver Bhojpuri data with the 13,000 gold HDTB sentences. Now I have a combined training corpus --- all in UD space.

\textit{Step 3.} Train a fresh Trankit model on this combined corpus. The Hindi HDTB gold acts as a regulariser that anchors the model to correct UD structure, while the silver Bhojpuri adds language-specific supervision.

Result: 54.27\% UAS and 36.70\% LAS on BHTB --- surpassing System A's 35.36\% LAS by 1.34 points. This is state-of-the-art on real Bhojpuri parsing.''
\end{pitchbox}

% ─────────────────────────────────────────────────
\section{Slide 16 --- Why UD-Bridge Works}

\begin{pitchbox}{What the slide shows}
The schema-consistency argument on the left; why silver data is good enough on the right.
\end{pitchbox}

\begin{pitchbox}{What to say}
``Let me explain why this approach works, because it's the main intellectual insight of Part 2.

System A is already a parser that lives in UD space, trained on Hindi HDTB. When it parses Bhojpuri, its output labels are automatically UD-compatible. Yes, they're noisy --- System A gets things wrong. But they're not \textit{systematically} wrong. They're just random errors scattered across a correct UD label inventory.

On the right, I address the obvious concern --- what makes this silver data good enough? Three reasons. The errors are random, not systematic. Adding HDTB gold pulls the silver errors back toward correct UD patterns --- it regularises the training. And crucially, because training and test are now in the same schema, any gain on the training objective transfers directly to BHTB evaluation.

The core claim, highlighted at the bottom: \textbf{noisy data in the right annotation space is more useful than high-quality data in the wrong annotation space} --- at least for benchmark evaluation.''
\end{pitchbox}

\begin{tipbox}
\textbf{Emphasise:} ``schema consistency'' is the methodological insight that generalises beyond Bhojpuri. State it clearly.
\end{tipbox}

% ─────────────────────────────────────────────────
\section{Slide 17 --- Results on BHTB Gold}

\begin{pitchbox}{What the slide shows}
A vs.\ K comparison table, why-K-wins bullets, and a headline win box.
\end{pitchbox}

\begin{pitchbox}{What to say}
``Here are the Part 2 results.

System A, zero-shot baseline: 52.78 UAS, 35.36 LAS.
System K, my UD-Bridge approach: 54.27 UAS, 36.70 LAS.

K wins by plus 1.49 UAS and plus 1.34 LAS.

Why K wins --- three reasons. One: it adds Bhojpuri-specific supervision from the silver data, which System A never sees. Two: training and test labels are both in the UD schema, so there's no mismatch. Three: HDTB gold anchors the training and stops silver noise from dominating.

The bottom line is in the red box: K beats A by 1.49 UAS and 1.34 LAS --- state-of-the-art on real Bhojpuri parsing.''
\end{pitchbox}

% ─────────────────────────────────────────────────
\section{Slide 18 --- Error Analysis: Where the Parser Wins and Loses}

\begin{pitchbox}{What the slide shows}
Per-relation LAS for System A, split into structural/local vs.\ clausal/long-range.
\end{pitchbox}

\begin{pitchbox}{What to say}
``To understand the parser's behaviour better, I broke down System A's errors by relation type. The table has two groups.

On the left, structural and local relations --- things like case markers, root assignment, punctuation, and noun modifiers. These transfer well from Hindi. Case markers at 83\%, root at 50\% --- quite good.

On the right, clausal and long-range relations --- relative clauses, adverbial clauses, complement clauses, coordination. These collapse. \textit{acl} at 1\%, \textit{advcl} at 4\%, \textit{xcomp} at 0\%, \textit{ccomp} at 0\%.

The pattern is clear. Short-arc, morphologically-marked relations transfer well because those features are consistent between Hindi and Bhojpuri. But clause-level relations depend on word order and syntactic structures where Bhojpuri diverges from Hindi --- and Hindi supervision alone can't teach that.

The lesson for future work: if we want the parser to improve on clausal relations, we need native Bhojpuri supervision --- which is exactly what UD-Bridge provides through its silver training data.''
\end{pitchbox}

% ─────────────────────────────────────────────────
\section{Slide 19 --- Part 2 Takeaway}

\begin{pitchbox}{What the slide shows}
A red takeaway box with Part 2 conclusion; what made the difference; bigger methodological lesson.
\end{pitchbox}

\begin{pitchbox}{What to say}
``Summarising Part 2: \textbf{UD-Bridge self-training beats the zero-shot baseline on BHTB --- advancing state-of-the-art on real Bhojpuri parsing.}

What made the difference? Three things: training and test schemas are aligned, silver Bhojpuri data provides language-specific supervision, and HDTB gold stops the silver noise from running away.

The bigger lesson --- and this is a methodological insight that applies beyond just Bhojpuri --- is that schema consistency matters as much as model quality. Noisy-but-correctly-labelled beats clean-but-mislabelled. And self-training is a principled, clean way to bridge zero-shot to low-resource settings.''
\end{pitchbox}

\clearpage

% ─────────────────────────────────────────────────
\begin{center}
{\LARGE\bfseries\color{accent}CONCLUSION}
\end{center}

% ─────────────────────────────────────────────────
\section{Slide 20 --- Key Findings Across Both Parts}

\begin{pitchbox}{What the slide shows}
Four numbered findings that cut across Parts 1 and 2.
\end{pitchbox}

\begin{pitchbox}{What to say}
``Stepping back across both parts, I have four main findings.

\textit{One: freezing the backbone is essential.} Full fine-tuning in System F collapsed to 17.75\% LAS because it destroys the multilingual prior. Adapters preserve it at negligible cost.

\textit{Two: syntax-aware alignment beats representation alignment.} SACT's Arc-KL and CTS supervise parsing decisions directly --- a stronger and faster signal than MSE on hidden states.

\textit{Three: schema consistency matters as much as model quality.} Strong performance on one dataset doesn't automatically transfer to a differently-labelled benchmark. The training signal must live in the test set's label space.

\textit{Four: self-training with schema-consistent silver data is a clean bridge into low-resource settings.} UD-Bridge shows that a zero-shot teacher can produce useful training data for a student, even without any gold Bhojpuri.''
\end{pitchbox}

% ─────────────────────────────────────────────────
\section{Slide 21 --- Contributions}

\begin{pitchbox}{What the slide shows}
Two architectural contributions (G, H), one benchmark contribution (K), and one methodological insight.
\end{pitchbox}

\begin{pitchbox}{What to say}
``My contributions group into three buckets.

From Part 1, two \textit{architectural} contributions. System G: my parallel adapter architecture with frozen XLM-R and MSE alignment. System H, the main one: SACT --- three syntax-aware alignment losses (content-word cosine, Arc-KL distillation, Cross-lingual Tree Supervision). H achieves 50.08\% LAS on aligned Dev.

From Part 2, one \textit{benchmark} contribution. System K: UD-Bridge silver self-training, achieving 36.70\% LAS on BHTB --- new state-of-the-art.

And across both parts, one \textit{methodological insight}: I identified and characterised the annotation-schema mismatch problem in cross-lingual parsing, and showed that noisy but schema-consistent supervision outperforms clean but schema-mismatched supervision on benchmarks. This generalises beyond just Bhojpuri.''
\end{pitchbox}

% ─────────────────────────────────────────────────
\section{Slide 22 --- Future Work}

\begin{pitchbox}{What the slide shows}
Four categories: short-term, data augmentation, architecture improvements, generalisation.
\end{pitchbox}

\begin{pitchbox}{What to say}
``Future work splits into four directions.

\textit{Short term:} iterative silver-to-model-to-silver self-training so the model bootstraps itself over multiple rounds. Agreement filtering between A and K predictions to identify reliable silver labels.

\textit{Data augmentation:} synthetic UD treebanks via CPG-to-UD conversion plus Hindi-to-Bhojpuri translation. And extending BHTB gold annotations.

\textit{Architecture improvements:} multi-layer adapters across all 12 XLM-R layers instead of just some. Non-projective MST decoding for more complex tree structures. Morphology-aware alignment.

\textit{Generalisation:} extend the framework to other low-resource Indic languages --- Maithili, Magahi, Awadhi, Rajasthani. Cluster adapters across language families. Multi-source transfer combining signals from several Indic languages.''
\end{pitchbox}

% ─────────────────────────────────────────────────
\section{Slide 23 --- Thank You}

\begin{pitchbox}{What the slide shows}
Closing slide with name, institute, supervisor, and repo name.
\end{pitchbox}

\begin{pitchbox}{What to say}
``That's the main talk. Thank you, sir. Happy to take any questions.''
\end{pitchbox}

\begin{tipbox}
\textbf{Emphasise:} short, confident closing. Don't trail off. Pause and make eye contact.
\end{tipbox}

\clearpage

% ─────────────────────────────────────────────────
\begin{center}
{\LARGE\bfseries\color{accent}BACKUP SLIDES}
\end{center}

\vspace{0.2cm}
\begin{tipbox}
Use backup slides only if specifically asked. Don't volunteer to show them. A one-line summary per slide is enough unless the supervisor wants more depth.
\end{tipbox}

% ─────────────────────────────────────────────────
\section{Slide 24 --- Backup: Hyperparameters}

\begin{pitchbox}{If asked}
``Adapter bottleneck dimension 64, standard Pfeiffer setup. Arc MLP dimension 500, label MLP dimension 100, MLP dropout 0.33. Max 60 epochs, batch size 16. Optimiser is AdamW with learning rate 2$\times$10\textsuperscript{-3} for adapters and 5$\times$10\textsuperscript{-5} for full fine-tuning. The LR choice reflects the fact that adapters are initialised near zero and can tolerate larger steps; full FT needs smaller steps to avoid breaking the pretrained backbone.''
\end{pitchbox}

% ─────────────────────────────────────────────────
\section{Slide 25 --- Backup: Datasets}

\begin{pitchbox}{If asked}
``HDTB --- Hindi, 13,304 sentences, used as training auxiliary in Part 1 and training main in Part 2. Hindi-Bhojpuri aligned corpus --- 30,966 pairs, Part 1 main training. Aligned Dev split --- 3,097 sentences, Part 1 evaluation. BHTB --- 357 sentences, UD, Part 2 evaluation, test-only. The key asymmetry is that only 357 gold Bhojpuri sentences exist in the UD universe --- which is exactly why transfer and self-training are essential.''
\end{pitchbox}

% ─────────────────────────────────────────────────
\section{Slide 26 --- Backup: SACT Loss in Detail}

\begin{pitchbox}{If asked}
``Three equations shown. Cosine loss is weighted cosine similarity between matched Hindi-Bhojpuri hidden states, with token-level weights. Arc-KL is KL divergence between Hindi teacher and Bhojpuri student arc score distributions, summed over tokens. CTS is the standard supervised cross-entropy loss on both arc and label, applied using the Hindi gold heads reused on Bhojpuri.

Final combined loss weights: Bhojpuri main 1.0, Hindi auxiliary 0.5, cosine 0.4, Arc-KL 2.0, CTS 0.2. Arc-KL has the highest weight because distillation provides the most direct supervision on parsing decisions --- even though its individual ablation drop is small at convergence, it dominates early training. The weights were selected on a small Dev grid.''
\end{pitchbox}

\clearpage

% ─────────────────────────────────────────────────
\section*{General Presenting Tips}

\begin{tcolorbox}[colback=softbg,colframe=accent,boxrule=0.3pt,title=\textbf{While speaking},coltitle=white,colbacktitle=accent]
\begin{itemize}[leftmargin=*,itemsep=3pt]
  \item Don't read the bullets verbatim. Speak the idea behind them.
  \item The narrative spine is \textbf{Part 1 $\to$ Part 2}. Finish Part 1 cleanly before starting Part 2.
  \item Slides that matter most: \textbf{Slide 8 (SACT)}, \textbf{Slide 12 (Part 1 Takeaway)}, \textbf{Slide 13 (BHTB schema mismatch)}, \textbf{Slide 15 (UD-Bridge)}. Slow down on these.
  \item On result slides (9, 17), lead with the headline number, \textit{then} the caveat.
  \item If interrupted, answer the question directly. Don't rush back to the bullets.
\end{itemize}
\end{tcolorbox}

\begin{tcolorbox}[colback=yellow!10,colframe=orange!70,boxrule=0.3pt,title=\textbf{Handling tough questions},coltitle=white,colbacktitle=orange!80!black]
\begin{itemize}[leftmargin=*,itemsep=3pt]
  \item \textit{``Why is F so weak?''} --- Small dataset plus 270M params plus no alignment signal. Overfitting and forgetting.
  \item \textit{``Is 1.34 LAS a significant gain on BHTB?''} --- Yes, because BHTB is only 357 sentences and the schema is strict. Small absolute LAS gains reflect real structural improvements. I can compute significance bootstraps if needed.
  \item \textit{``Why not just annotate more Bhojpuri gold?''} --- We propose exactly that in Future Work. But UD-Bridge works with what's available today --- that's its practical value.
  \item \textit{``Does this transfer to Maithili/Magahi/etc?''} --- Expected to, because the schema-consistency principle is language-agnostic. That's part of future work.
\end{itemize}
\end{tcolorbox}

\begin{tcolorbox}[colback=part1color!5,colframe=part1color,boxrule=0.3pt,title=\textbf{One-sentence elevator pitch (if asked to compress)},coltitle=white,colbacktitle=part1color]
\textit{``I built a cross-lingual dependency parser for Bhojpuri by combining syntax-aware adapter alignment (SACT) on a Hindi-Bhojpuri aligned corpus with UD-Bridge silver self-training on the BHTB benchmark, achieving 50.08 LAS on aligned Dev and 36.70 LAS on BHTB --- state-of-the-art for Bhojpuri.''}
\end{tcolorbox}

\end{document}
